\section{Turing Machine --- Build --- Seam}
%%% generated by the blueprint pipeline from TCSlib.Complexity.TuringMachine.Build.Seam
\subsection*{Turing Machine --- Build --- Seam: Introduction}

Declarations of the Lean module \texttt{TCSlib.Complexity.TuringMachine.Build.Seam}.

\subsection*{Turing Machine --- Build --- Seam: Results}

\begin{definition}[Machine composition across an exit state]
\lean{Turing.seamCompTM}
\label{Turing.seamCompTM}
\leanok
\uses{Turing.MultiTapeTM}
Let $M_1$ and $M_2$ be deterministic Turing machines with the same input alphabet and
number of work tapes, with state sets $S_1$ and $S_2$. Given an exit state $e\in S_1$
and an entry state $b\in S_2$, form a machine with state set $S_1\sqcup S_2$ that starts
in the initial state of $M_1$. In a state from $S_1$ other than $e$, it performs the
transition of $M_1$, retaining its input movement, work-tape actions, and output, and
places any successor state in the $S_1$ summand. At $e$, it leaves the input and work
tapes in place, emits nothing, and enters $b$ in the $S_2$ summand; in the $S_2$ summand
it performs $M_2$'s transitions with successor states placed in that summand.
\end{definition}

\begin{lemma}[Step in the Left State Component]
\lean{Turing.seamComp_step_left}
\label{Turing.seamComp_step_left}
\leanok
\uses{Turing.Cfg, Turing.Cfg.mapState, Turing.MultiTapeTM, Turing.MultiTapeTM.step, Turing.seamCompTM}
\difficulty{4}
Let $M_1$ and $M_2$ be machines with $k$ work tapes over a common alphabet, with state
sets $S_1,S_2$. Choose an exit state $exit\in S_1$ and entry state $entry\in S_2$. For
every configuration $c$ of $M_1$ whose state is not $exit$, one step of the
seam-composed machine from $c$ with its state embedded in the left component equals the
configuration obtained by taking one step of $M_1$ from $c$ and embedding the resulting
state in the left component.
\end{lemma}

\begin{lemma}[Step in the Right State Component]
\lean{Turing.seamComp_step_right}
\label{Turing.seamComp_step_right}
\leanok
\uses{Turing.Cfg, Turing.Cfg.mapState, Turing.MultiTapeTM, Turing.MultiTapeTM.step, Turing.seamCompTM}
\difficulty{4}
Let $M_1$ and $M_2$ be machines with $k$ work tapes over a common alphabet, with state
sets $S_1,S_2$. Choose an exit state $exit\in S_1$ and entry state $entry\in S_2$. For
every configuration $c$ of $M_2$, one step of the seam-composed machine from $c$ with
its state embedded in the right component equals the configuration obtained by taking
one step of $M_2$ from $c$ and embedding the resulting state in the right component.
\end{lemma}

\begin{lemma}[Stationary Action on a Configuration]
\lean{Turing.seam_stationary_apply}
\label{Turing.seam_stationary_apply}
\leanok
\uses{Turing.Action.apply, Turing.Cfg}
\difficulty{1}
Let $c$ be a configuration and $q$ a state. Applying the action that makes no input-head
movement, makes no work-tape writes or head movements, emits no output symbol, and sets
the successor state to $q$ yields $c$ with its state replaced by $q$.
\end{lemma}

\begin{lemma}[Seam Handoff at the Exit State]
\lean{Turing.seamComp_dispatch}
\label{Turing.seamComp_dispatch}
\leanok
\uses{Turing.Cfg, Turing.Cfg.mapState, Turing.MultiTapeTM, Turing.MultiTapeTM.step, Turing.seamCompTM, Turing.seam_stationary_apply}
\difficulty{3}
Let $M_1$ and $M_2$ be machines with $k$ work tapes over a common alphabet, with state
sets $S_1,S_2$. Choose an exit state $exit\in S_1$ and entry state $entry\in S_2$. For
every configuration $c$ of $M_1$ whose state is $exit$, one step of the seam-composed
machine from $c$ with its state embedded in the left component equals the configuration
obtained by replacing $c$'s state by $entry$ and embedding it in the right component.
\end{lemma}

\begin{theorem}[Composed run agrees with the first machine]
\lean{Turing.seamCompTM_runFrom_left}
\label{Turing.seamCompTM_runFrom_left}
\leanok
\uses{Turing.Cfg, Turing.Cfg.mapState, Turing.MultiTapeTM, Turing.MultiTapeTM.runFrom, Turing.MultiTapeTM.runFrom_succ_eq_step', Turing.seamCompTM, Turing.seamComp_step_left}
\difficulty{4}
Let $M_1$ and $M_2$ be deterministic machines with $k$ work tapes and a common input
alphabet, with state sets $S_1$ and $S_2$. Form a machine on $S_1\sqcup S_2$ that
follows $M_1$ on left states, switches from the specified exit state to the specified
entry state in one step without changing the tapes, and follows $M_2$ on right states.
If a run of $M_1$ from configuration $c$ avoids the exit state at every time $t<T$, then
for every $t\le T$ the composed machine, started from $c$ with its state placed in the
left component, is exactly the configuration reached by $M_1$ at time $t$, with its
state placed in the left component.
\end{theorem}

\begin{theorem}[Composed run after the transition]
\lean{Turing.seamCompTM_runFrom_right}
\label{Turing.seamCompTM_runFrom_right}
\leanok
\uses{Turing.Cfg, Turing.Cfg.mapState, Turing.MultiTapeTM, Turing.MultiTapeTM.runFrom, Turing.MultiTapeTM.runFrom_add, Turing.MultiTapeTM.runFrom_comm_of_step, Turing.MultiTapeTM.runFrom_succ_eq_step', Turing.seamCompTM, Turing.seamCompTM_runFrom_left, Turing.seamComp_dispatch, Turing.seamComp_step_right}
\difficulty{3}
Let $M_1$ and $M_2$ be deterministic machines with $k$ work tapes and a common input
alphabet, with state sets $S_1$ and $S_2$. Form a machine on $S_1\sqcup S_2$ that
follows $M_1$ on left states, switches from the specified exit state to the specified
entry state in one step without changing the tapes, and follows $M_2$ on right states.
If $M_1$ reaches configuration $c_1$ from $c_0$ after $T_1$ steps, first visiting the
exit state at time $T_1$, then for every natural number $t$ the composed machine started
from $c_0$ with its state placed in the left component, after $T_1+1+t$ steps, equals
the configuration reached by $M_2$ after $t$ steps from $c_1$ with its state replaced by
the entry state, with that resulting state placed in the right component.
\end{theorem}

\begin{lemma}[Run Through a Seam Handoff]
\lean{Turing.seamComp_run_general}
\label{Turing.seamComp_run_general}
\leanok
\uses{Turing.Cfg, Turing.Cfg.mapState, Turing.MultiTapeTM, Turing.MultiTapeTM.runFrom, Turing.seamCompTM, Turing.seamCompTM_runFrom_right}
\difficulty{1}
Let $M_1$ and $M_2$ be machines with $k$ work tapes over a common alphabet, with state
sets $S_1,S_2$. Choose $exit\in S_1$ and $entry\in S_2$. Suppose the run of $M_1$ from
configuration $c_0$ reaches $c_1$ after $T_1$ steps, with state $exit$, and avoids
$exit$ at every earlier time. Suppose the run of $M_2$ from $c_1$ with its state
replaced by $entry$ reaches configuration $c_3$ after $T_2$ steps. Then the
seam-composed machine, started from $c_0$ with its state embedded in the left component,
reaches $c_3$ with its state embedded in the right component after $T_1+1+T_2$ steps.
\end{lemma}

\begin{lemma}[No Early Visit to the Second Exit State]
\lean{Turing.seamComp_firstReturn_general}
\label{Turing.seamComp_firstReturn_general}
\leanok
\uses{Turing.Cfg, Turing.Cfg.mapState, Turing.MultiTapeTM, Turing.MultiTapeTM.runFrom, Turing.seamCompTM, Turing.seamCompTM_runFrom_left, Turing.seamCompTM_runFrom_right}
\difficulty{5}
Let $M_1$ and $M_2$ be machines with $k$ work tapes over a common alphabet, with state
sets $S_1,S_2$. Choose $exit\in S_1$ and $entry,q_2\in S_2$. Suppose the run of $M_1$
from configuration $c_0$ reaches configuration $c_1$ at time $T_1$, with state $exit$,
and avoids $exit$ at every time $t<T_1$. Suppose also that the run of $M_2$ from $c_1$
with its state replaced by $entry$ avoids $q_2$ at every time $t<T_2$. Then the
seam-composed run from $c_0$ with state embedded in the left component does not have
state $q_2$ embedded in the right component at any time $t<T_1+1+T_2$.
\end{lemma}

\begin{lemma}[Visited Tape Cells Under Seam Composition]
\lean{Turing.seamComp_visited_general}
\label{Turing.seamComp_visited_general}
\leanok
\uses{Turing.Cfg, Turing.Cfg.mapState, Turing.MultiTapeTM, Turing.MultiTapeTM.runFrom, Turing.MultiTapeTM.visitedByTapeHead, Turing.seamCompTM, Turing.seamCompTM_runFrom_left, Turing.seamCompTM_runFrom_right}
\difficulty{4}
Let $M_1$ and $M_2$ be machines with the same work tapes, and let $e$ be a state of
$M_1$ and $a$ a state of $M_2$. Suppose a computation of $M_1$ from $c_0$ reaches $c_1$
after $T_1$ steps, with state $e$, and does not visit state $e$ at any earlier time
$t<T_1$. For each work tape $i$, the set of head positions visited by the seam-composed
machine from the state-renamed configuration $c_0$ through time $T_1+1+T_2$ is contained
in the union of the positions visited by $M_1$ from $c_0$ through time $T_1$ and those
visited by $M_2$ from the state-renamed configuration with state $a$ and the other data
of $c_1$ through time $T_2$.
\end{lemma}

\begin{lemma}[State Renaming of a Word Configuration]
\lean{Turing.seam_ofWords_mapState}
\label{Turing.seam_ofWords_mapState}
\leanok
\uses{Turing.Cfg.mapState, Turing.Cfg.ofWords}
\difficulty{1}
Let $f$ map a state set $S_1$ to a state set $S_3$, let $q\in S_1$, and let $w$ be a
family of $k$ Boolean words. Renaming the state of the word configuration for $q$ and
$w$ by $f$ gives the word configuration for $f(q)$ and $w$.
\end{lemma}

\begin{theorem}[Run of a composed machine]
\lean{Turing.seamCompTM_run}
\label{Turing.seamCompTM_run}
\leanok
\uses{Turing.Cfg.mapState, Turing.Cfg.ofWords, Turing.MultiTapeTM, Turing.MultiTapeTM.runFrom, Turing.seamCompTM, Turing.seamComp_run_general, Turing.seam_ofWords_mapState}
\difficulty{3}
Let $M_1$ and $M_2$ be deterministic machines with $k$ work tapes over the Boolean
alphabet, with state sets $S_1$ and $S_2$, and let $x$ be a Boolean input word. Form a
machine on $S_1\sqcup S_2$ that follows $M_1$ on left states, switches from the
specified exit state to the specified entry state in one step without changing the
tapes, and follows $M_2$ on right states. If $M_1$ takes the initial configuration with
state $start$ and work tapes $w_0$ to the configuration with state $exit$ and work tapes
$w_1$ after $T_1$ steps, never having state $exit$ earlier, and $M_2$ takes the initial
configuration with state $entry$ and work tapes $w_1$ to the configuration with state
$q_2$ and work tapes $w_2$ after $T_2$ steps, then the composed machine takes the
initial configuration with state $(left,start)$ and work tapes $w_0$ to the
configuration with state $(right,q_2)$ and work tapes $w_2$ after $T_1+1+T_2$ steps.
\end{theorem}

\begin{theorem}[First return in sequential machine composition]
\lean{Turing.seamCompTM_firstReturn}
\label{Turing.seamCompTM_firstReturn}
\leanok
\uses{Turing.Cfg.mapState, Turing.Cfg.ofWords, Turing.MultiTapeTM, Turing.MultiTapeTM.runFrom, Turing.seamCompTM, Turing.seamComp_firstReturn_general, Turing.seam_ofWords_mapState}
\difficulty{3}
Let $k\in\mathbb N$, let $x$ be a Boolean input word, and let $M_1,M_2$ be deterministic
$k$-work-tape machines over Boolean symbols, with state sets $S_1,S_2$ and a decidable
equality on $S_1$. Choose an exit state $e\in S_1$, an entry state $u\in S_2$, initial
state $s\in S_1$, final state $q\in S_2$, three families of $k$ Boolean work-tape words
$w_0,w_1,w_2$, and times $T_1,T_2\in\mathbb N$. Suppose $M_1$, started in state $s$ with
tape words $w_0$, reaches the configuration with state $e$ and words $w_1$ at time
$T_1$, and does not have state $e$ at any earlier time. Suppose also that $M_2$, started
in state $u$ with words $w_1$, reaches state $q$ and words $w_2$ at time $T_2$, and does
not have state $q$ at any earlier time. Then the machine formed by switching from $M_1$
at state $e$ to $M_2$ at state $u$ does not reach state $q$ in its $S_2$ component at
any time $t<T_1+1+T_2$ when started from state $s$ with words $w_0$.
\end{theorem}

\begin{theorem}[Visited Work Tape Cells Under Seam Composition]
\lean{Turing.seamCompTM_visitedByTapeHead}
\label{Turing.seamCompTM_visitedByTapeHead}
\leanok
\uses{Turing.Cfg.ofWords, Turing.MultiTapeTM, Turing.MultiTapeTM.runFrom, Turing.MultiTapeTM.visitedByTapeHead, Turing.seamCompTM, Turing.seamComp_visited_general, Turing.seam_ofWords_mapState}
\difficulty{1}
Let $M_1$ and $M_2$ be machines with $k$ work tapes and Boolean input $x$, with states
in $S_1$ and $S_2$, respectively. Suppose that starting $M_1$ in state $start$ with
work-tape words $w_0$ reaches state $exit$ with words $w_1$ after $T_1$ steps, and that
$M_1$ is not in state $exit$ at any earlier time $t<T_1$. Suppose also that starting
$M_2$ in state $entry$ with words $w_1$ reaches state $q_2$ with words $w_2$ after $T_2$
steps. Then, for every work tape $i$, the cells visited by its head in the seam-composed
machine through time $T_1+1+T_2$ are contained in the union of the cells visited by that
tape in the first computation through time $T_1$ and in the second computation through
time $T_2$.
\end{theorem}

\begin{theorem}[Space on one tape under composition]
\lean{Turing.seamCompTM_spaceUsedByTape_le_add}
\label{Turing.seamCompTM_spaceUsedByTape_le_add}
\leanok
\uses{Turing.Cfg.ofWords, Turing.MultiTapeTM, Turing.MultiTapeTM.runFrom, Turing.MultiTapeTM.spaceUsedByTape, Turing.seamCompTM, Turing.seamCompTM_visitedByTapeHead}
\difficulty{1}
Let $M_1$ and $M_2$ be deterministic machines with $k$ work tapes over the Boolean
alphabet, with state sets $S_1$ and $S_2$, and let $x$ be a Boolean input word. Form a
machine on $S_1\sqcup S_2$ that follows $M_1$ on left states, switches from the
specified exit state to the specified entry state in one step without changing the
tapes, and follows $M_2$ on right states. Suppose $M_1$ takes the initial configuration
with state $start$ and work tapes $w_0$ to $(exit,w_1)$ after $T_1$ steps, without
visiting $exit$ earlier, and $M_2$ takes $(entry,w_1)$ to $(q_2,w_2)$ after $T_2$ steps.
For each work tape $i$, the number of distinct cells visited by its head in the composed
run through time $T_1+1+T_2$ is at most the sum of the numbers of distinct cells visited
in the two respective runs through times $T_1$ and $T_2$.
\end{theorem}

\begin{theorem}[Total space under machine composition]
\lean{Turing.seamCompTM_spaceUsed_le_add}
\label{Turing.seamCompTM_spaceUsed_le_add}
\leanok
\uses{Turing.Cfg.ofWords, Turing.MultiTapeTM, Turing.MultiTapeTM.runFrom, Turing.MultiTapeTM.spaceUsed, Turing.seamCompTM, Turing.seamCompTM_spaceUsedByTape_le_add}
\difficulty{3}
Let $M_1$ and $M_2$ be deterministic machines with $k$ work tapes over the Boolean
alphabet, with state sets $S_1$ and $S_2$, and let $x$ be a Boolean input word. Form a
machine on $S_1\sqcup S_2$ that follows $M_1$ on left states, switches from the
specified exit state to the specified entry state in one step without changing the
tapes, and follows $M_2$ on right states. Suppose $M_1$ takes $(start,w_0)$ to
$(exit,w_1)$ after $T_1$ steps, without visiting $exit$ earlier, and $M_2$ takes
$(entry,w_1)$ to $(q_2,w_2)$ after $T_2$ steps. Then the total number of distinct
work-tape cells visited, summed over all $k$ tapes, in the composed run through time
$T_1+1+T_2$ is at most the sum of the total visited-cell counts in the two runs through
times $T_1$ and $T_2$, respectively.
\end{theorem}

\begin{theorem}[Maximum space on one tape under composition]
\lean{Turing.seamCompTM_spaceUsedByTape_le_max}
\label{Turing.seamCompTM_spaceUsedByTape_le_max}
\leanok
\uses{Turing.Cfg.ofWords, Turing.MultiTapeTM, Turing.MultiTapeTM.runFrom, Turing.MultiTapeTM.spaceUsedByTape, Turing.MultiTapeTM.visitedByTapeHead, Turing.seamCompTM, Turing.seamCompTM_visitedByTapeHead}
\difficulty{4}
Let $M_1$ and $M_2$ be deterministic machines with $k$ work tapes over the Boolean
alphabet, with state sets $S_1$ and $S_2$, and let $x$ be a Boolean input word. Form a
machine on $S_1\sqcup S_2$ that follows $M_1$ on left states, switches from the
specified exit state to the specified entry state in one step without changing the
tapes, and follows $M_2$ on right states. Suppose $M_1$ takes $(start,w_0)$ to
$(exit,w_1)$ after $T_1$ steps, without visiting $exit$ earlier, and $M_2$ takes
$(entry,w_1)$ to $(q_2,w_2)$ after $T_2$ steps. For each work tape $i$, if either the
set of cells visited by its head in the first run through time $T_1$ or the set visited
in the second run through time $T_2$ is exactly the singleton consisting of cell $0$,
then the composed run through time $T_1+1+T_2$ visits at most the maximum of the two
runs' visited-cell counts on tape $i$.
\end{theorem}

\begin{theorem}[Composition of two machine runs]
\lean{Turing.seamCompTM_run_ofCfg}
\label{Turing.seamCompTM_run_ofCfg}
\leanok
\uses{Turing.Cfg, Turing.Cfg.mapState, Turing.MultiTapeTM, Turing.MultiTapeTM.runFrom, Turing.seamCompTM, Turing.seamComp_run_general}
\difficulty{1}
Let $M_1$ and $M_2$ be deterministic machines with $k$ work tapes and a common input
alphabet, with state sets $S_1$ and $S_2$. Form a machine on $S_1\sqcup S_2$ that
follows $M_1$ on left states, switches from the specified exit state to the specified
entry state in one step without changing the tapes, and follows $M_2$ on right states.
If $M_1$ reaches configuration $c_1$ from $c_0$ after $T_1$ steps, first visiting the
exit state at time $T_1$, and $M_2$ reaches $c_3$ after $T_2$ steps from $c_1$ with its
state replaced by the entry state, then the composed machine, started from $c_0$ with
its state placed in the left component, reaches $c_3$ with its state placed in the right
component after $T_1+1+T_2$ steps.
\end{theorem}

\begin{theorem}[First return during machine composition]
\lean{Turing.seamCompTM_firstReturn_ofCfg}
\label{Turing.seamCompTM_firstReturn_ofCfg}
\leanok
\uses{Turing.Cfg, Turing.Cfg.mapState, Turing.MultiTapeTM, Turing.MultiTapeTM.runFrom, Turing.seamCompTM, Turing.seamComp_firstReturn_general}
\difficulty{1}
Let $M_1$ and $M_2$ be deterministic machines with $k$ work tapes and a common input
alphabet, with states in sets $S_1$ and $S_2$. Form a machine with state set $S_1\sqcup
S_2$ that follows $M_1$ on left states, switches from the exit state of $M_1$ to the
specified entry state of $M_2$ in one step without changing the tapes, and follows $M_2$
on right states. If $M_1$ reaches configuration $c_1$ from $c_0$ after $T_1$ steps,
first visits the exit state at that time, and $M_2$, started from the tapes and head
positions of $c_1$ with its state replaced by the entry state, reaches $c_3$ after $T_2$
steps with state $q_2$, and does not visit $q_2$ earlier, then the composed machine does
not have state $q_2$ in its right component at any time $t<T_1+1+T_2$ when started from
$c_0$ with its state placed in the left component.
\end{theorem}

\begin{theorem}[Visited Cells for General Configurations]
\lean{Turing.seamCompTM_visitedByTapeHead_ofCfg}
\label{Turing.seamCompTM_visitedByTapeHead_ofCfg}
\leanok
\uses{Turing.Cfg, Turing.Cfg.mapState, Turing.MultiTapeTM, Turing.MultiTapeTM.runFrom, Turing.MultiTapeTM.visitedByTapeHead, Turing.seamCompTM, Turing.seamComp_visited_general}
\difficulty{1}
Let $M_1$ and $M_2$ be machines with $k$ work tapes over the same alphabet, with state
sets $S_1$ and $S_2$, and choose states $exit\in S_1$ and $entry\in S_2$. Let $c_0,c_1$
be configurations of $M_1$ on the same input, and suppose $M_1$ reaches $c_1$ from $c_0$
after $T_1$ steps, with state $exit$, while avoiding state $exit$ at every time $t<T_1$.
For every natural number $T_2$ and work-tape index $i$, the cells visited by tape $i$ in
the seam-composed run from $c_0$ with its states embedded in the left component, through
time $T_1+1+T_2$, are contained in the union of the cells visited by $M_1$ from $c_0$
through time $T_1$ and those visited by $M_2$ from $c_1$ with its state replaced by
$entry$, through time $T_2$.
\end{theorem}

\begin{definition}[Two-seam machine composition]
\lean{Turing.seamCompTM₂}
\leanok
\uses{Turing.MultiTapeTM, Turing.seamCompTM}
Let $M_0,M_1,M_2$ be deterministic Turing machines with a common input alphabet and the
same number of work tapes, with state sets $S_0,S_1,S_2$. Choose exit states $e_1,e_2\in
S_0$ and entry states $b_1\in S_1,b_2\in S_2$. The resulting machine, with states
$(S_0\sqcup S_1)\sqcup S_2$, first composes $M_0$ with $M_1$ by sending state $e_1$ to
$b_1$, then composes that machine with $M_2$ by sending the state corresponding to $e_2$
in the $S_0$ summand to $b_2$.
\end{definition}

\begin{theorem}[Two-Stage Seam Composition Run]
\lean{Turing.seamCompTM₂_run_ofCfg}
\leanok
\uses{Turing.Cfg, Turing.Cfg.mapState, Turing.MultiTapeTM, Turing.MultiTapeTM.runFrom, Turing.seamCompTM, Turing.seamCompTM_runFrom_left, Turing.seamCompTM_run_ofCfg, Turing.seamCompTM₂}
\difficulty{6}
Let $M_0,M_1,M_2$ be machines with $k$ work tapes over a common alphabet and respective
state sets $S_0,S_1,S_2$. Choose states $e_1,e_2\in S_0$ with $e_1\ne e_2$, and entry
states $entry_1\in S_1$ and $entry_2\in S_2$. If a run of $M_0$ from configuration $c_0$
reaches $c_1$ after $T_0$ steps and avoids both states $e_1$ and $e_2$ at every earlier
time, then the following two conclusions hold: if $c_1$ has state $e_1$, the twice
seam-composed machine, run from $c_0$ with its state embedded in the leftmost component,
agrees after the handoff and any $t$ further steps with the run of $M_1$ from $c_1$ with
state replaced by $entry_1$, embedded in the middle component; if $c_1$ has state $e_2$,
it agrees instead with the run of $M_2$ from $c_1$ with state replaced by $entry_2$,
embedded in the rightmost component. In each case the claimed agreement holds for every
natural number $t$, at time $T_0+1+t$.
\end{theorem}

\begin{definition}[Machine with a Released Anchor State]
\lean{Turing.seamReleaseTM}
\label{Turing.seamReleaseTM}
\leanok
\uses{Turing.MultiTapeTM}
Given a deterministic Turing machine $M$ with $k$ work tapes and state set $S$, and a
chosen state $a\in S$, there is a machine with state set $\{\ast\}\sqcup S$ whose
initial state is $\ast$. In state $\ast$ it uses the transition of $M$ at $a$, and in
state $s\in S$ it uses the transition of $M$ at $s$; in both cases it preserves the
input-head movement, work-tape actions, and output, and any successor state is placed in
the $S$ summand.
\end{definition}

\begin{theorem}[Seam-Release Step from the Fresh State]
\lean{Turing.seamReleaseTM_step_fresh}
\label{Turing.seamReleaseTM_step_fresh}
\leanok
\uses{Turing.Cfg, Turing.Cfg.mapState, Turing.MultiTapeTM, Turing.MultiTapeTM.step, Turing.seamReleaseTM}
\difficulty{3}
Let $M$ be a machine, let $a$ be one of its states, and let $c$ be a configuration for
$M$ in state $a$. One step of the seam-release machine from the configuration obtained
by renaming $c$'s state into the left summand gives the configuration obtained by taking
one step of $M$ from $c$ and renaming the resulting state into the right summand.
\end{theorem}

\begin{theorem}[Seam-Release Step in the Original-State Summand]
\lean{Turing.seamReleaseTM_step_inr}
\label{Turing.seamReleaseTM_step_inr}
\leanok
\uses{Turing.Cfg, Turing.Cfg.mapState, Turing.MultiTapeTM, Turing.MultiTapeTM.step, Turing.seamReleaseTM}
\difficulty{4}
Let $M$ be a machine with state set $S$, let $a$ be a state of $M$, and let $c$ be any
configuration for $M$. One step of the seam-release machine from the configuration
obtained by renaming $c$'s state into the right summand equals the configuration
obtained by taking one step of $M$ from $c$ and renaming its state into the right
summand.
\end{theorem}

\begin{theorem}[Seam-Release Machine Run]
\lean{Turing.seamReleaseTM_runFrom_pos}
\label{Turing.seamReleaseTM_runFrom_pos}
\leanok
\uses{Turing.Cfg, Turing.Cfg.mapState, Turing.MultiTapeTM, Turing.MultiTapeTM.runFrom, Turing.MultiTapeTM.runFrom_comm_of_step, Turing.MultiTapeTM.runFrom_succ_eq_step, Turing.seamReleaseTM, Turing.seamReleaseTM_step_fresh, Turing.seamReleaseTM_step_inr}
\difficulty{4}
Let $M$ be a multi-tape Turing machine with state set $S$, let $a\in S$, and let $c$ be
a configuration of $M$ on an input word $x$ whose state is $a$. For every positive
integer $t$, running the seam-release machine associated to $M$ and $a$ for $t$ steps
from the configuration obtained by placing the state of $c$ in the left summand yields
the configuration obtained by running $M$ from $c$ for $t$ steps and placing its state
in the right summand.
\end{theorem}

\begin{theorem}[First Return in Seam Release]
\lean{Turing.seamReleaseTM_firstReturn}
\label{Turing.seamReleaseTM_firstReturn}
\leanok
\uses{Turing.Cfg, Turing.Cfg.mapState, Turing.MultiTapeTM, Turing.MultiTapeTM.runFrom, Turing.seamReleaseTM, Turing.seamReleaseTM_runFrom_pos}
\difficulty{5}
Let $M$ be a machine, let $a$ be one of its states, and suppose a computation from $c$
starts in state $a$, runs for a positive number $T$ of steps to a configuration $c'$ in
state $a$, and has no intermediate return to state $a$ at any time $t$ with $0<t<T$.
Starting from the configuration obtained by renaming the state of $c$ to the left
summand, the seam-release machine reaches the state-renamed version of $c'$ in the right
summand after $T$ steps; at every time $t<T$, its state is not the right-summand copy of
$a$.
\end{theorem}

\begin{theorem}[Tape-Head Visits Under Seam Release]
\lean{Turing.seamReleaseTM_visitedByTapeHead}
\label{Turing.seamReleaseTM_visitedByTapeHead}
\leanok
\uses{Turing.Cfg, Turing.Cfg.mapState, Turing.MultiTapeTM, Turing.MultiTapeTM.visitedByTapeHead, Turing.seamReleaseTM, Turing.seamReleaseTM_runFrom_pos}
\difficulty{4}
Let $M$ be a machine, let $a$ be one of its states, and let $c$ be a configuration for
$M$ in state $a$. For every time $t$ and every work tape $i$, the set of head positions
visited through time $t$ by the seam-release machine started from $c$ with its state
renamed into the left summand is exactly the set visited by $M$ started from $c$.
\end{theorem}

\begin{definition}[Phase-controlled machine]
\lean{Turing.phaseHostTM}
\label{Turing.phaseHostTM}
\leanok
\uses{Turing.Action.mapState, Turing.MultiTapeTM}
Let $I$ be a set of phases, let each phase $i\in I$ have a state set $\mathrm{Ph}(i)$
and a $k$-work-tape Turing machine over the alphabet $\Gamma$, and let $q_0$ be an
initial state in the disjoint union of these phase state sets. Let a rule assign to each
state in phase $i$ either no phase change or a state in some phase. The resulting
machine starts at $q_0$: when the rule gives a state, it moves to that state with the
input head and every work head staying still, writes nothing, and emits nothing; when
the rule gives no state, it performs the corresponding phase machine’s transition, with
any successor state kept in the same phase.
\end{definition}

\begin{theorem}[Runs of a phase-hosted machine]
\lean{Turing.phaseHostTM_run}
\label{Turing.phaseHostTM_run}
\leanok
\uses{Turing.Cfg, Turing.Cfg.mapState, Turing.MultiTapeTM, Turing.MultiTapeTM.runFrom, Turing.MultiTapeTM.runFrom_mapState_of_agreeOn, Turing.MultiTapeTM.runFrom_succ_eq_step', Turing.MultiTapeTM.step, Turing.phaseHostTM}
\difficulty{4}
Let $k$ be a natural number, let $\Gamma$ and $I$ be types, let $\mathrm{Ph}(i)$ be a
state type for each $i\in I$, let $M_i$ be a $k$-work-tape machine with state type
$\mathrm{Ph}(i)$, and let $q_0$ be a state tagged by an element of $I$. Suppose $c,c'$
are configurations of $M_i$ on input $x$, $T\in\mathbb{N}$, and running $M_i$ from $c$
for $T$ steps yields $c'$. If $s'$ is the state of $c'$ and the reduction rule sends
$(i,s')$ to $q'$, and if at every time $u<T$ the run of $M_i$ from $c$ is in a state $s$
for which the reduction rule sends $(i,s)$ to no state, then for every $u\le T$ the
phase-hosted machine's run from the configuration $c$ tagged by $i$ agrees with the run
of $M_i$ from $c$ tagged by $i$, and after $T+1$ steps it is $c'$ with its state
replaced by $q'$.
\end{theorem}
